\begin{lemma}[Change lemma, mixed location: largest closed starred stage]
Fix a finite loopless multigraph with its exact edge set, a palette of
$k$ colors, and a distinguished edge $e_0$.  Let $(T,\varphi_n)$ be a
positive-rung extended Tashkinov tree.  Let $\mathcal B$ be a hierarchy
base and let
\[
 T_n=T_{n,0}\subset T_{n,1}\subset\cdots\subset
 T_{n,q}\subset T_{n,q+1}=T
\]
be a good hierarchy under $\varphi_n$, where $q$ is its last divider
index and
\[
 A=T_{n,q}^{*}=\noderef{HierarchyStarredStage}(H,q).
\]
Assume that $T$ has maximum property, that $V(T)$ is not elementary,
that $\varphi_n$ is proper and leaves exactly $e_0$ uncolored, and that
the rung number is positive.  These data form a
\noderef{HierarchyCounterexample} $C$.

Assume the six lower-rung assertions of the technical theorem, uniformly
for every positive-rung maximum ETT on the same graph, palette, and
distinguished edge, with arbitrary certified input and generating
colorings and with rung number smaller than that of $T$: lower-rung
elementarity, terminal interchangeability, the two PE stability and
stored-pair assertions, the terminal PE component assertion, the
defective-endpoint order assertion, and the mod-coloring assertion.
Assume also the same-rung induction hypothesis that every maximum ETT
with a good hierarchy having fewer than $q$ dividers is elementary.
Finally assume the large-color estimate
\[
 2n+11\leq |\overline\varphi_n(T_n)|.
\]

Put $\ell=|T|-|A|$, the number of exact tail additions above $A$, and let
\[
 T(y_{\ell-1})=T\big|_{\min\{|A|+(\ell-1),|T|\}}
\]
be the exact ordered prefix supplied by \noderef{TreeSequencePrefix}.
Assume (6.6): $V(T(y_{\ell-1}))$ is elementary under $\varphi_n$.
Assume moreover that $C$ is lexicographically minimal among all
\noderef{HierarchyCounterexample}s on the fixed graph, palette, and
$e_0$ whose ETTs have the same rung number as $T$ and which have exactly
$q$ dividers, allowing their certified input colorings, generating
colorings, bases, ETTs, and good hierarchies to vary.  Explicitly, no
such same-rung counterexample $C'$ satisfies
\[
 \noderef{HierarchyCounterexampleLess}
  \bigl(\noderef{HierarchyCounterexampleMeasure}(C'),
        \noderef{HierarchyCounterexampleMeasure}(C)\bigr).
\]

Let $\alpha\ne\beta$ be colors missing on $T(y_{\ell-1})$.  Write
\[
 i_c=\noderef{HierarchyFirstMissingPosition}(\varphi_n,T,c),
 \qquad v_c=T.\mathit{vertex}(i_c).
\]
If at least one of $\alpha,\beta$ is not missing on $A$, assume both
\[
 v_\alpha\prec v_\beta                                      \tag{a}
\]
and the exact edge-identity exclusion
\[
 \alpha\notin
 \left\{c:\exists t\in\operatorname{Fin}|T|,
       |A|\leq t<i_\beta,
       \ \varphi_n(T.\mathit{edge}(t))=c\right\}.             \tag{b}
\]
Thus (b) is precisely
$\alpha\notin\varphi_n\langle T(v_\beta)-A\rangle$; it is not
an endpoint-only abbreviation.


Assume in addition that the first $\alpha$-missing vertex lies in the
last starred stage and the first $\beta$-missing vertex lies outside it:
\[
        v_\alpha\in V(A),\qquad v_\beta\notin V(A).
\]
Assume that $q>0$.  Let $r$ satisfy $0\leq r<q$, assume that
$\{\alpha\}$ is closed on $T_{n,r}^{*}$, and assume explicitly that
for every $s$ with $r<s\leq q$, the singleton $\{\alpha\}$ is not
closed on $T_{n,s}^{*}$.  Thus $r$ is the largest closed starred-stage
index; no missingness at stage $r$ is assumed.  Then
\[
 \noderef{TwoColorComponent}(G,\varphi_n,\alpha,\beta,
       \operatorname{inl}v_\alpha)
 =
 \noderef{TwoColorComponent}(G,\varphi_n,\alpha,\beta,
       \operatorname{inl}v_\beta).
\]
\end{lemma}
\begin{proof}
Write $H=C.\mathit{hierarchy}$ and $q=H.\mathit{dividerCount}$, and
retain the notation $A,\ell,T(y_{\ell-1}),i_\alpha,i_\beta,
v_\alpha,v_\beta$ from the statement.  Put
\[
 K_x=\noderef{TwoColorComponent}
       (G,\varphi_n,\alpha,\beta,\operatorname{inl}x),
 \qquad S_r=\noderef{HierarchyStarredStage}(H,r).
\]
Thus $S_r=T_{n,r}^{*}$.  All component supports below are supports in
the incidence graph, so their edge-nodes retain individual parallel-edge
identities.

We first record the exact prefix and path facts.  Since $q>0$,
\noderef{HierarchyStarredStage} identifies $A$ with $T_{n,q}$.
The nestedness fields of \noderef{TashkinovHierarchy} make every
ordinary and starred stage through $q$ an exact ordered initial segment
of $A$, and make $A$ a proper initial segment of $T$.  Hence
$\ell=|T|-|A|>0$.  The definition of the exact prefix in
\noderef{TreeSequencePrefix}, together with
\noderef{TreeSequenceInitialSegment}, puts $A$ and every earlier stage
inside $T(y_{\ell-1})$ entry by entry.  Restricting the assumed
\noderef{ColoredVertexSetElementary} property makes every such stage
elementary under $\varphi_n$.

The nonempty-candidate case of
\noderef{HierarchyFirstMissingPosition} says that $v_\alpha$ misses
$\alpha$, $v_\beta$ misses $\beta$, and no earlier displayed vertex
misses the corresponding color.  Their mixed locations make them
distinct.  Elementarity of $T(y_{\ell-1})$ says that $v_\alpha$ does
not miss $\beta$ and $v_\beta$ does not miss $\alpha$.  Expanding
\noderef{MissingColors} and \noderef{PresentColors}, there is therefore
a $\beta$-colored edge incident with $v_\alpha$ and an $\alpha$-colored
edge incident with $v_\beta$, while no $\alpha$-edge is incident with
$v_\alpha$ and no $\beta$-edge is incident with $v_\beta$.
Properness makes the displayed incident edge unique at each endpoint.
Equivalently, each endpoint vertex-node has degree exactly one in its
two-color incidence component; the global upper bound is also the
conclusion of \noderef{TwoColorIncidenceDegreeAtMostTwo}.  Thus each
component satisfies the leaf condition in
\noderef{TwoColorPathComponent}.  The displayed incident edge-node
makes it edge-bearing, so both $K_{v_\alpha}$ and $K_{v_\beta}$ satisfy
\noderef{TwoColorNontrivialPathComponent}.  The spanning paths supplied
by \noderef{TwoColorComponentSpanningPath} therefore represent exactly
the two paper path components to which
\noderef{TreeInterchangeable} will be applied.

We next construct the predecessor data needed below.  Write $n>0$ for
the rung number and $\varphi_{n-1}$ for the coloring in the slot preceding
the displayed generating coloring $\varphi_n$.  Eliminate the last
positive certificate in \noderef{TashkinovSeries}; by
\noderef{TashkinovSeriesStep} it is exactly one of SE, PE, or RE, and
the generating-coloring field of \noderef{ExtendedTashkinovTree}
identifies its current coloring with $\varphi_n$.

If $n>1$, restrict every array of the certified series to slots through
$n-1$, retaining its tree slots through $T_n$.  Every retained
certificate has the same constructor and the same witnesses because it
only refers to retained slots.  The last retained noninitial certificate
gives the strict exact inclusion from $T_{n-1}$ to $T_n$.  Use $T_n$
itself as the displayed terminal tree; the upper inclusion in its last
certified closure is reflexive.  These data give an actual
\noderef{ExtendedTashkinovTree} $\mathcal E^{-}$ of rung $n-1$, exact
terminal tree $T_n$, and generating coloring $\varphi_{n-1}$.  The
initial maximum-order equality in \noderef{MaximumProperty} is unchanged,
and every closure bound indexed below $n-1$ is literally the corresponding
old bound.  Hence every field of maximum property restricts to
$\mathcal E^{-}$.  Lower assertion (i), applied to this actual maximum
ETT, gives
\[
             V(T_n)\text{ is elementary under }\varphi_{n-1}. \tag{1}
\]
If $n=1$, slot zero of the given series is its certified initial tuple
and its closure is exactly $T_1=T_n$.  The elementarity conclusion of
\noderef{TechnicalTheoremBase}, instantiated with this very initial
tuple, gives (1).  No zero-rung ETT or maximum-property assertion at
rung zero is introduced.

The color $\alpha$ is missing on $A=T_{n,q}$ because its witness
$v_\alpha$ lies there.  Apply \noderef{rutcor} using (1), the full-tail
elementarity hypothesis (6.6), the stated good hierarchy, and the actual
maximal closed index $r$.  Its first conclusion gives the exact
edge-identity exclusion
\[
 \alpha\notin
   \varphi_n\langle T_{n,q}-S_r\rangle.                         \tag{2}
\]

We now derive the missingness needed for interchangeability; singleton
closedness alone is not used as a missingness assertion.  We prove by
descending induction that $\alpha$ is missing on every $T_{n,s}^{*}$
for $r\leq s\leq q$.  At $s=q$, the witness is the actual vertex
$v_\alpha\in A$.  Suppose $r<s\leq q$ and $\alpha$ is missing on
$T_{n,s}^{*}=T_{n,s}$ but not on $T_{n,s-1}^{*}$.  A vertex witnessing
the first fact lies in the exact vertex difference
$V(T_{n,s})-V(T_{n,s-1}^{*})$.  The stage $T_{n,s}$ is elementary, so
\noderef{GoodHierarchyNewMissingClosed} makes $\alpha$ closed on
$T_{n,s}=T_{n,s}^{*}$.  This contradicts the stated nonclosedness at
every index strictly larger than $r$.  Descending to $s=r+1$ proves
\[
             \alpha\in\overline\varphi_n(S_r).                 \tag{3}
\]
Choose a vertex of $S_r$ witnessing (3).  The least-candidate property
of \noderef{HierarchyFirstMissingPosition}, together with the entrywise
prefix relation, places $v_\alpha$ itself in $S_r$.

We prove that $\alpha,\beta$ are $S_r$-interchangeable.  If $r>0$,
then $S_r=T_{n,r}$, and \noderef{step1} applies with (3), the assumed
singleton closedness, the restricted elementarity, maximum property,
the six lower-rung assertions, the fewer-divider assertion, and the
large-color estimate.

Suppose $r=0$.  If the actual terminal constructor is PE,
\noderef{HierarchyBase} identifies $S_r$ with the exact core
$T_n\vee R_n$.  Reflexivity gives both stable relations in
\noderef{CombinedStableColoring} at the identity coloring
$\varphi_n$.  Conclusion (iii) of \noderef{interchange}, applied with
(3), the assumed singleton closedness, maximum property, and the stated
lower package, gives interchangeability on this core.

It remains to check the actual SE and RE constructors at $r=0$.  Their
nonparallel base equation gives $S_r=T_n$.  When $n>1$, use the maximum
ETT $\mathcal E^{-}$ constructed above.  In SE, the last constructor's
\noderef{MaximumDefectiveChoice} identifies its selected coloring with
$\varphi_n$; its `selectedProper', `selectedOneUncolored', and
`selectedStable' fields say respectively that $\varphi_n$ is proper,
leaves exactly $e_0$ uncolored, and is
$(T_n,D_{n-1},\varphi_{n-1})$-stable.  In RE, the constructor gives the
literal equality $\varphi_n=\varphi_{n-1}$, and reflexivity in
\noderef{sc2} supplies that same stable relation, while the retained
constructor supplies properness and one-uncoloredness.  Lower assertion
(ii), applied to $\mathcal E^{-}$ and this exact stable coloring, now
gives $T_n$-interchangeability because (3) makes $\alpha$ missing there.
When $n=1$, use the same SE fields or RE equality as the proper,
one-uncolored, and stable inputs of the interchangeability conclusion of
\noderef{TechnicalTheoremBase} for the actual initial closure $T_1=T_n$.
Consequently, in every case,
\[
        \alpha,\beta\text{ are }S_r\text{-interchangeable}.     \tag{4}
\]

Suppose for contradiction that $K_{v_\alpha}\ne K_{v_\beta}$.
The first certified path component meets $S_r$ at $v_\alpha$.  If the
second met $S_r$, (4) and \noderef{TreeInterchangeable} would identify
the two components.  It therefore contains no vertex-node of $S_r$.
By \noderef{TwoColorComponentsDisjoint} it is disjoint from the first
component as well.  Moreover, it contains no edge-node incident with
$S_r$: such an edge-node is adjacent in the two-color incidence graph to
its endpoint vertex-node in $S_r$, which would put that vertex-node in
the same component.

Define
\[
              \sigma_n=\varphi_n/K_{v_\beta}
\]
using \noderef{TwoColorComponentSwap}.  The theorem
\noderef{TwoColorComponentSwapProper} makes $\sigma_n$ proper.  The
switch fixes the value \textsf{none} and exchanges only the two nonempty
colors, so expanding \noderef{IsOneEdgeUncolored} shows that exactly
$e_0$ remains uncolored.  The disjoint component through $v_\alpha$ is
fixed, so $v_\alpha$ still misses $\alpha$; at $v_\beta$ the switch
replaces missing $\beta$ by missing $\alpha$.

Let $B=T(v_\beta)$ be the exact ordered prefix ending at the displayed
vertex $v_\beta$.  Since $v_\beta\in T(y_{\ell-1})-A$, if
$s_B=|B|-|A|$, then
\[
                         1\leq s_B\leq\ell-1<\ell.             \tag{5}
\]
Thus $A\prec B\preccurlyeq T$.  Both distinguished vertices lie in
$B$ and miss $\alpha$ under $\sigma_n$, so $B$ is not elementary.

Before changing any ETT or hierarchy data, apply
\noderef{GoodHierarchyTerminalPrefixRepackage} to this exact prefix.
It gives an old-color ETT $\mathcal E_B$ on the actual tree $B$, with
the same complete certified series and maximum property, a copied base
$\mathcal B_B$, a same-divider hierarchy $H_B$ whose nonterminal and
starred stages are the old ones, and a genuine old-color good witness
$\mathcal G_B$.  It also gives the old computed stored-set,
stage-missing-set, and Gamma equalities.  Its segment conclusion is the
fallback-sensitive containment needed for the old truncation, not an
equality of full \noderef{HierarchySegmentColors} functions.

We next establish the exact stability data.  The terminal ladder tree
$T_n$ and, in PE, both $T_n$ and $R_n$ lie in the core and hence in
$S_r$.  The switched full component is disjoint from each of these
trees and has no edge-node incident with them.  Applying
\noderef{OutsideKempeRecoloringStable} to $T_n$, and separately to
$R_n$ in PE, checks both the vertexwise missing-set clause and the
controlled incident-edge clause of \noderef{StableColoring}.  Thus
\[
 \sigma_n\text{ is }(T_n,D_n,\varphi_n)\text{-stable},          \tag{6}
\]
and in PE the two relations together give
\noderef{CombinedStableColoring}.  In that case
\noderef{CombinedStableClosurePreserved} also fixes every displayed
identity of the two constituents and transports $R_n$ as the same exact
closure under $\sigma_n$.

The conditional ordered-unused hypothesis of the statement applies
because $v_\beta\notin A$ implies that $\beta$ is missing nowhere on
$A$.  It excludes $\alpha$ from the exact interval $B-A$.  Joining that
interval to (2), whose endpoint is $A=T_{n,q}$, gives
\[
 \alpha\notin\varphi_n\langle B-S_r\rangle.                    \tag{7}
\]

We now construct the recolored augmentation route, rather than infer it
from the final vertex set.  Let $T_{n+1}$ be the final closure stored by
the actual last constructor isolated above.  In PE its exact seed is
$P=T_n$; in SE or RE it is the exact one-edge seed $P=T_n+f_n$.
Unfolding that constructor's `closedExtension' through
\noderef{TashkinovClosure} first supplies a fixed-$\varphi_n$
\noderef{TashkinovAugmentationSequence} from $P$ to $T_{n+1}$.
The seed is an exact prefix of $B$.  In PE this follows from
$T_n\preccurlyeq S_r\preccurlyeq A\prec B$.  In SE/RE, both the
one-edge seed and $B$ are exact prefixes of $T_{n+1}$, and
$|T_n+f_n|=|T_n|+1\leq|A|<|B|$ because $q>0$; entrywise comparison in
their common extension using \noderef{TreeSequenceInitialSegment} gives
$P\preccurlyeq B$.  The `belowFinal' field of
\noderef{ExtendedTashkinovTree} gives
$B\preccurlyeq T\preccurlyeq T_{n+1}$.  Hence
\noderef{TashkinovAugmentationSequenceRestrict}, still at the old fixed
coloring, gives a $\varphi_n$-augmentation sequence
\[
        P=X_0\longrightarrow X_1\longrightarrow\cdots
          \longrightarrow X_m=B.                              \tag{8}
\]

We first compare the colorings on the seed.  Every displayed identity of
$T_n$ is fixed by (6).  The only additional identity in an SE/RE seed is
$f_n$.  If its position is below $|S_r|$, component disjointness fixes it.
The only case not covered that way is $r=0$, where the nonparallel base
gives $S_r=T_n$.  Equation (7) excludes
$\varphi_n(f_n)=\alpha$.  If the value lies outside
$\{\alpha,\beta\}$, \noderef{TwoColorComponentSwap} fixes it.  Suppose
its value is $\beta$.  The first-missing location of $v_\beta$ outside
$T_n$ says that $\beta$ is not missing on $T_n$.  In SE,
\noderef{MaximumDefectiveChoice} gives
$\varphi_n(f_n)=\delta_n$ and the constructor puts
$S_n=\{\delta_n\}$.  In RE, the
\noderef{RevisitingExtensionCertificate} puts $f_n$ on its recorded
two-color cycle and the constructor stores both cycle colors.  In either
case, unfolding \noderef{TashkinovStoredColors} shows
$\beta\in D_n$.  The same constructor makes the exact identity $f_n$
incident with $T_n$, so the incident-edge clause of the stable relation
(6) gives $\sigma_n(f_n)=\varphi_n(f_n)=\beta$.  This checks the paper's
SE/RE first-seed exception and proves equality on all of $P$.

Induct through the reflexive-transitive chain (8).  For one old step
$X_i\longrightarrow X_{i+1}$, unfold
\noderef{TashkinovAugmentation}.  It supplies the exact new last edge
$g$, its old color $c=\varphi_n(g)$, and a vertex
$w\in V(X_i)$ at which $c$ is missing under $\varphi_n$.
If the new position is below $|S_r|$, then $X_i\preccurlyeq S_r$;
component disjointness fixes $g$ and every incident edge at $w$ whose
color is $\alpha$ or $\beta$, while every other color is globally fixed.
Thus $\sigma_n(g)=c$ and $c$ is still missing at $w$.

If the new position is at least $|S_r|$, it lies strictly before the
terminal position of $v_\beta$ in $B$.  Equation (7) gives
$c\ne\alpha$.  If $c=\beta$, the old augmentation witness $w$ occurs in
the preceding exact prefix and would be a vertex before $v_\beta$ missing
$\beta$, contrary to the least-candidate clause of
\noderef{HierarchyFirstMissingPosition}.  Hence
$c\notin\{\alpha,\beta\}$.  The component swap fixes $g$ and every
incident $c$-edge at $w$, so once again $\sigma_n(g)=c$ and $c$ remains
missing there.  In both positional cases the length, attachment, and
ordered-prefix clauses are coloring-independent.  Therefore the same
exact step is a $\sigma_n$-\noderef{TashkinovAugmentation}.  Induction
proves a fixed-$\sigma_n$ augmentation sequence from the appropriate
seed $P$ to $B$.  Simultaneously, the seed check and the new-edge check
at each indexed step prove
\[
 \sigma_n(g)=\varphi_n(g)
 \quad\text{for every displayed edge identity }g\text{ of }B.  \tag{9}
\]
In particular, every old-prefix missing witness required by the route
was checked at the step where it is used.

Apply \noderef{Td} to the maximum ETT $\mathcal E_B$ using (6), the
stated lower package, properness, and one-uncoloredness.  It gives a
\noderef{ModColoring}.  If the terminal constructor is PE, apply
\noderef{HierarchySplitter} to $\mathcal E_B$, its copied base and
hierarchy, the verified combined stability, and the $\sigma_n$ route
from $T_n$ to $B$.  It returns an actual
\noderef{CorrespondingExtendedTashkinovTree} $\mathcal E'_B$ on $B$,
maximum property, a base $\mathcal B'_B$, and a hierarchy $K$ satisfying
\noderef{SameHierarchyStages} with $H_B$.

If the constructor is SE or RE, use in the universal non-PE clause of
\noderef{ModColoring} the exact identity $f_n$, the constructor's
equality $F_n=\{f_n\}$, the exact seed $T_n+f_n$, and the route just
proved.  It returns a corresponding ETT $\mathcal E'_B$ on $B$.
\noderef{MaximumPropertyTransport} gives maximum property.  Construct
$\mathcal B'_B$ with core equal to its replacement terminal ladder tree,
copy the unused closure and support position, use terminal stability for
the missing-set compatibility, and use the retained non-PE kind for the
nonparallel clause.  Copy the divider count, all nonterminal stages,
terminal $B$, and all divider positions from $H_B$ to form $K$; all
nestedness and exact divider-prefix fields copy, and both cores are the
same $T_n$.  These data satisfy
\noderef{SameHierarchyStages}.  Thus both terminal-kind branches have
produced the same replacement package without applying the PE splitter
to a non-PE constructor.

It remains to construct goodness under $\sigma_n$.  First compare all
sets used as indices.  The `storedColorsAgree' clauses of
\noderef{CorrespondingExtendedTashkinovTree} identify every raw certified
$S_i$ set.  Terminal stability identifies missing colors on $T_n$;
unfolding \noderef{TashkinovStoredColors}, namely raw stored union minus
the current $T_n$-missing set, therefore identifies the old and new
terminal $D_n$ sets.

Every ordinary or starred hierarchy stage lies in $A$.  At a stage
vertex outside $K_{v_\beta}$, all incident values are fixed.  At a stage
vertex inside that component, $\beta$ was present because its first
missing vertex is the later vertex $v_\beta$, and $\alpha$ was present
because the elementary prefix already has the distinct
$\alpha$-missing vertex $v_\alpha$ in the other component.  Properness
then says that the complete component switch exchanges one incident edge
of each affected color, leaving both colors present; all other colors are
fixed.  Hence the vertex missing-color set is unchanged at every vertex
of every ordinary and starred stage.  In particular, their
\noderef{TreeMissingColors} sets agree.

Now unfold \noderef{HierarchyStoredColorsAt}: it subtracts the missing
set of the relevant starred stage from the just-identified terminal
$D_n$.  Thus, for every $0\leq j\leq q$,
\[
       D^K_{n,j}=D^{H_B}_{n,j}.                                \tag{10}
\]
The same vertexwise computation on $T_n$ at index zero and on each
positive ordinary stage, followed by
\noderef{HierarchyStageMissingColors}, identifies all stage-missing
sets.  These equalities are obtained before any Gamma family is copied.
Transport the Gamma function of $\mathcal G_B$ along the divider-count
equality and use the same pair at every common index (10).

We next compute the totalized first-missing positions in $B$.  First,
no vertex before $v_\alpha$ misses $\alpha$ under $\varphi_n$.  If such
a vertex became $\alpha$-missing under $\sigma_n$, it would lie in the
switched component and would have been $\beta$-missing under
$\varphi_n$, contradicting that $v_\beta$ is the first such vertex and
occurs after $v_\alpha$.  Outside the component nothing changed, while
$v_\alpha$ lies in the disjoint component and still misses $\alpha$.
Thus $v_\alpha$ remains the first $\alpha$-missing vertex in $B$.

Second, no vertex of $B$ misses $\beta$ under $\sigma_n$.  A vertex
outside the switched component is fixed and occurs no later than the old
first $\beta$-missing vertex.  If a vertex in the switched component
became $\beta$-missing, it was $\alpha$-missing under $\varphi_n$.
Elementarity of $T(y_{\ell-1})$ forces it to equal $v_\alpha$, but that
vertex lies in the other component.  At $v_\beta$ itself, the old
$\alpha$-edge whose existence was proved above becomes a $\beta$-edge.
The no-candidate clause of \noderef{HierarchyFirstMissingPosition}
therefore assigns $\beta$ the last position of $B$, exactly the old
displayed position of $v_\beta$.  Finally, for every
$\eta\notin\{\alpha,\beta\}$, the component swap fixes every incident
$\eta$-edge globally.  Its missing candidates in $B$ are pointwise
identical under the two colorings, so their least positions agree when
nonempty and their last-position fallbacks agree when empty.

We verify each field of \noderef{GoodHierarchy}.  For avoidance, fix
$j\leq q$, a relevant color $\eta\in D^K_{n,j}$, and one exact edge
position witnessing a color in the new
\noderef{HierarchySegmentColors}(K,j,\eta)$.  If $\eta=\alpha$, the
first-position calculation identifies its upper bound with the old
position of $v_\alpha$; if $\eta=\beta$, it identifies the new fallback
with the old terminal position $v_\beta$; for every other $\eta$, the
candidate calculation identifies the two bounds.  Same hierarchy stages
identify the starred and next-stage bounds and the exact edge at this
position, and (9) identifies its color.  Thus this individual edge and
color witness membership in the corresponding old segment.  The old
avoidance field of $\mathcal G_B$, the stage-missing equality, and the
copied pair prove new avoidance.  This is the needed pointwise
new-to-old containment at relevant indices, and does not assert equality
of the full segment-color functions.

The `twoColors' and `disjointPairs' fields copy literally from
$\mathcal G_B$ using (10).  For `uniqueRefresh', same stages identify the
exact stage differences, (10) identifies their stored sets, the
vertexwise calculation identifies the missing-color unions on those
differences, and the Gamma pairs agree.  Hence the old unique witness and
its uniqueness proof copy at every positive index.

For the PE `parallelBase' cardinality field, the component is disjoint
from the core $T_n\vee R_n\preccurlyeq S_r$.  The constituent missing
sets are therefore unchanged; (10) and Gamma equality identify the
reserved union, so the old inequality is the new one.  For its
non-minus-closed field, unfold the old witness from $\mathcal G_B$ using
\noderef{TreeColorSetMinusClosed}, \noderef{TreeColorSetClosed}, and
\noderef{BoundaryColorEdges}.  It supplies a color $c$ missing on the
core and outside the old reserved union, and an individual exact boundary
edge $g$ of the core with $\varphi_n(g)=c$.  By
\noderef{BoundaryEdges} and \noderef{TreeIncidentEdge}, this same identity
is incident with a core vertex.  Component disjointness from the core
therefore says its edge-node is not switched, so
$\sigma_n(g)=\varphi_n(g)=c$.  Core missingness, the reserved union, and
the graph-theoretic boundary are unchanged.  The same $c,g$ consequently
witness failure of minus-closedness under $\sigma_n$.  In SE/RE the PE
premise of `parallelBase' is false.

Finally fix a positive stage $T_{n,j}$.  If $j\leq r$, it lies inside
$S_r$.  Every boundary identity incident with it is fixed by component
disjointness, while its missing set and reserved union have already been
identified.  Expanding \noderef{TreeColorSetMinusClosed} and
\noderef{TreeColorSetClosed}, any new counterexample color and exact
boundary edge would therefore be an old counterexample to the old
`minusClosed' field.

Suppose $j>r$.  The descending argument above says that $\alpha$ is
missing on $T_{n,j}$, while the branch hypothesis says that its singleton
is not closed there.  The old minus-closed field, unfolded through the
same two definitions, forces
\[
 \alpha\in
   \bigcup_{\eta\in D_{n,j}}\Gamma^{j-1}_\eta.                 \tag{11}
\]
Also $\beta$ is missing nowhere on $T_{n,j}\preccurlyeq A$, because its
first missing vertex is outside $A$.  To prove the new minus-closed
field, take a color $c$ missing under $\sigma_n$ on this stage and outside
the common reserved union, and suppose an exact boundary identity $g$
has $\sigma_n(g)=c$.  Stage missing-set equality makes $c$ old-missing.
If $c=\alpha$, (11) contradicts that it is outside the reserved union.
If $c=\beta$, the established nonmissingness of $\beta$ contradicts its
membership in the common stage-missing set.  Otherwise
$c\notin\{\alpha,\beta\}$, so \noderef{TwoColorComponentSwap} fixes the
exact identity $g$ and gives $\varphi_n(g)=c$; the same
\noderef{BoundaryColorEdges} witness contradicts the old minus-closed
field.  These cases prove `minusClosed' for every $j>r$.

We have now checked `twoColors', `avoidance', `disjointPairs',
`uniqueRefresh', both `parallelBase' clauses, and `minusClosed'.  They
construct a complete good witness $\mathcal G'_B$ under $\sigma_n$.
Only now apply \noderef{SameGoodHierarchy}: same hierarchy stages,
(10), and the transported Gamma equations give exactly its pointwise
data.

The resulting ETT on $B$, its maximum property, base, hierarchy, newly
constructed good witness, and the distinct vertices $v_\alpha,v_\beta$
both missing $\alpha$ form a \noderef{HierarchyCounterexample} $C'$ with
the same divider count $q$.  The explicit rounds equality returned by
\noderef{GoodHierarchyTerminalPrefixRepackage} identifies the rung of
$\mathcal E_B$ with that of the original ETT, and the `rungAgreement'
field of the \noderef{CorrespondingExtendedTashkinovTree} certificate
identifies the rung of $\mathcal E'_B$ with that of $\mathcal E_B$.
Consequently the ETT underlying $C'$ has the same rung as the ETT
underlying $C$.  It remains to compare the actual measures.
Put $p=\noderef{HierarchyPathNumber}(H)$ and let $p'$ be the path number
of $K$.  If $s_B<p$, the fallback bound in the path-number definition
gives $p'\leq s_B<p$.  If $p\leq s_B$, the old path tail beginning at
$p$ restricts to the exact prefix $B$, with its singleton endpoint still
counted as a path, so $p'\leq p$.  If equality holds, the second
coordinate of \noderef{HierarchyCounterexampleMeasure} strictly drops:
since $q>0$ gives $A=T_{n,q}$, equation (5) yields
\[
       |B|-|T_{n,q}|=s_B<\ell=|T|-|T_{n,q}|.
\]
Therefore \noderef{HierarchyCounterexampleLess} holds from $C'$ to $C$.
The same-rung equality just proved and the divider-count equality make
$C'$ eligible for the stated minimality, a contradiction.  The
supposition that the two path
components were distinct was false, and hence
$K_{v_\alpha}=K_{v_\beta}$ as required.
\end{proof}
